%============================================================
%  TalentLink – A Professional Freelancer Platform
%  Industrial Oriented Mini Project Report
%  JYOTHISHMATHI INSTITUTE OF TECHNOLOGY AND SCIENCE
%  Department of CSE | 2025-2026
%  Formatted as per JITS Mini-Project Guidelines
%============================================================

\documentclass[12pt,a4paper,oneside]{report}

%------------------------------------------------------------
% PACKAGES
%------------------------------------------------------------
\usepackage[left=1.5in, right=1.0in, top=1.0in, bottom=1.0in]{geometry}
\usepackage{times}           % Times New Roman font
\usepackage{fontenc}
\usepackage[T1]{fontenc}
\usepackage[utf8]{inputenc}
\usepackage{setspace}        % Line spacing
\usepackage{titlesec}        % Custom section/chapter titles
\usepackage{tocloft}         % TOC / LOF / LOT customisation
\usepackage{fancyhdr}        % Header / footer control
\usepackage{graphicx}        % Include images
\usepackage{float}           % Force figure placement
\usepackage{caption}         % Caption formatting
\usepackage{tabularx}        % Tables
\usepackage{booktabs}        % Professional table lines
\usepackage{array}           % Extended column definitions
\usepackage{enumitem}        % List customisation
\usepackage{hyperref}        % Clickable links and TOC
\usepackage{parskip}         % Space between paragraphs
\usepackage{indentfirst}     % Indent first paragraph of each section
\usepackage{microtype}       % Better text justification
\usepackage{url}             % URL formatting
\usepackage{longtable}       % Multi-page tables
\usepackage{amsmath}         % Math support
\usepackage{xcolor}          % Color support
\usepackage{emptypage}       % Suppress headers/footers on empty pages
\usepackage{afterpage}       % Insert commands after current page

%------------------------------------------------------------
% HYPERREF SETUP
%------------------------------------------------------------
\hypersetup{
    colorlinks=true,
    linkcolor=black,
    urlcolor=black,
    citecolor=black,
    pdfborder={0 0 0}
}

%------------------------------------------------------------
% FONT SIZE DEFINITIONS (per JITS guideline)
%  Chapter names  : 16pt Bold ALL CAPS
%  Main headings  : 14pt Bold ALL CAPS
%  Subheadings    : 14pt Bold Title Case
%  Sub-subheadings: 12pt Bold Title Case
%  Body text      : 12pt
%  Fig/Table cap  : 10pt Italic Title Case
%------------------------------------------------------------

%------------------------------------------------------------
% CHAPTER TITLE FORMAT  (16pt TNR Bold, centred)
%------------------------------------------------------------
\titleformat{\chapter}[block]
  {\centering\fontsize{16}{20}\selectfont\bfseries}  % format
  {Chapter~\thechapter}                               % label
  {0pt}                                               % sep
  {\vspace{6pt}\centering\MakeUppercase}              % before-code
  []                                                  % after-code

\titlespacing*{\chapter}{0pt}{0pt}{2\baselineskip}

%------------------------------------------------------------
% SECTION FORMAT  (14pt TNR Bold ALL CAPS)
%------------------------------------------------------------
\titleformat{\section}
  {\fontsize{14}{18}\selectfont\bfseries\MakeUppercase}
  {\thesection}{1em}{}

\titlespacing*{\section}{0pt}{1.5\baselineskip}{0.5\baselineskip}

%------------------------------------------------------------
% SUBSECTION FORMAT  (14pt TNR Bold Title Case)
%------------------------------------------------------------
\titleformat{\subsection}
  {\fontsize{14}{18}\selectfont\bfseries}
  {\thesubsection}{1em}{}

\titlespacing*{\subsection}{0pt}{1.2\baselineskip}{0.4\baselineskip}

%------------------------------------------------------------
% SUBSUBSECTION FORMAT  (12pt TNR Bold Title Case)
%------------------------------------------------------------
\titleformat{\subsubsection}
  {\fontsize{12}{16}\selectfont\bfseries}
  {\thesubsubsection}{1em}{}

\titlespacing*{\subsubsection}{0pt}{1\baselineskip}{0.3\baselineskip}

%------------------------------------------------------------
% CAPTION FORMAT  (10pt Italic Title Case)
%------------------------------------------------------------
\captionsetup{
    font={small,it},
    labelfont={small,it},
    justification=centering,
    skip=6pt
}

%------------------------------------------------------------
% PARAGRAPH SETTINGS
%------------------------------------------------------------
\setlength{\parindent}{1em}    % Single tab indent for new paragraphs
\setlength{\parskip}{6pt}       % Space between paragraphs

%------------------------------------------------------------
% SPACING
%------------------------------------------------------------
\onehalfspacing          % 1.5 line spacing for body text

%------------------------------------------------------------
% TOC / LOF / LOT FORMATTING
%------------------------------------------------------------
\renewcommand{\contentsname}{\centering\fontsize{16}{20}\selectfont\bfseries TABLE OF CONTENTS}
\renewcommand{\listfigurename}{\centering\fontsize{16}{20}\selectfont\bfseries LIST OF FIGURES}
\renewcommand{\listtablename}{\centering\fontsize{16}{20}\selectfont\bfseries LIST OF TABLES}

%------------------------------------------------------------
% FIGURE NUMBERING: Chapter.Figure  e.g. Fig 5.1
%------------------------------------------------------------
\numberwithin{figure}{chapter}
\numberwithin{table}{chapter}

%------------------------------------------------------------
% PAGE STYLE
%------------------------------------------------------------
\fancypagestyle{plain}{
  \fancyhf{}
  \fancyfoot[C]{\thepage}
  \renewcommand{\headrulewidth}{0pt}
}

\pagestyle{fancy}
\fancyhf{}
\fancyfoot[C]{\thepage}
\renewcommand{\headrulewidth}{0pt}

%============================================================
\begin{document}
%============================================================

%------------------------------------------------------------
%  FRONT MATTER  (Roman numerals, no page number on
%  Title / Certificate / Acknowledgement pages)
%------------------------------------------------------------
\pagenumbering{roman}
\setcounter{page}{1}

%============================================================
%  TITLE PAGE
%============================================================
\thispagestyle{empty}
\begin{titlepage}
  \centering
  \vspace*{0.3cm}

  % College Logo placeholder (replace with actual logo file)
  % \includegraphics[width=3cm]{logo.png}\\[0.3cm]

  {\fontsize{16}{20}\selectfont\bfseries
    JYOTHISHMATHI INSTITUTE OF TECHNOLOGY AND SCIENCE}\\[0.1cm]
  {\fontsize{12}{16}\selectfont
    (Autonomous, NBA (CSE, ECE, EEE) and NAAC `A' Grade)\\
    (Approved by AICTE, New Delhi, Affiliated to JNTUH, Hyderabad)\\
    Nustulapur, Karimnagar -- 505481, Telangana, India}\\[0.5cm]

  {\fontsize{14}{18}\selectfont\bfseries
    DEPARTMENT OF COMPUTER SCIENCE AND ENGINEERING}\\[0.8cm]

  \hrule height 1.5pt
  \vspace{0.3cm}
  {\fontsize{14}{18}\selectfont\bfseries\itshape
    Industrial Oriented Mini Project Report}\\[0.1cm]
  {\fontsize{14}{18}\selectfont on}\\[0.3cm]
  {\fontsize{16}{20}\selectfont\bfseries
    TALENT LINK -- A PROFESSIONAL FREELANCER PLATFORM}\\[0.3cm]
  \hrule height 1.5pt
  \vspace{0.6cm}

  {\fontsize{12}{16}\selectfont
    Submitted to\\[0.2cm]
    {\bfseries Jawaharlal Nehru Technological University, Hyderabad}\\[0.2cm]
    For the partial fulfilment of requirements for the award of the degree in\\[0.2cm]
    {\bfseries BACHELOR OF TECHNOLOGY}\\[0.1cm]
    in\\[0.1cm]
    {\bfseries COMPUTER SCIENCE AND ENGINEERING}}\\[0.8cm]

  \begin{minipage}[t]{0.45\textwidth}
    \centering
    {\fontsize{12}{16}\selectfont\bfseries Submitted by}\\[0.4cm]
    {\fontsize{12}{16}\selectfont
      ADI ROSHINI \hfill (23271A05D9)\\[0.2cm]
      VATARIKARI MANISHA \hfill (23271A05C8)\\[0.2cm]
      ARMOOR AKHILA \hfill (23271A05B9)\\[0.2cm]
      ARANDKAR VAESHNAVI \hfill (23271A05G3)}
  \end{minipage}
  \hfill
  \begin{minipage}[t]{0.45\textwidth}
    \centering
    {\fontsize{12}{16}\selectfont\bfseries Under the Esteemed Guidance of}\\[0.4cm]
    {\fontsize{12}{16}\selectfont
      \textbf{DR. M. RAVINDAR}\\
      Associate Professor\\
      Dept. of CSE}
  \end{minipage}

  \vfill
  {\fontsize{12}{16}\selectfont\bfseries 2025--2026}
\end{titlepage}

%============================================================
%  CERTIFICATE
%============================================================
\thispagestyle{empty}
\newpage
\begin{center}
  {\fontsize{16}{20}\selectfont\bfseries CERTIFICATE}
\end{center}
\vspace{2\baselineskip}

\begin{doublespace}
{\fontsize{12}{16}\selectfont
This is to certify that the Industrial Oriented Mini Project Report entitled
\textbf{``Talent Link -- A Professional Freelancer Platform''} is being submitted by
\textbf{ADI ROSHINI} (23271A05D9), \textbf{VATARIKARI MANISHA} (23271A05C8),
\textbf{ARMOOR AKHILA} (23271A05B9), \textbf{ARANDKAR VAESHNAVI} (23271A05G3),
in partial fulfilment of the requirements for the award of the Degree of
\textbf{Bachelor of Technology} in \textbf{Computer Science \& Engineering} to the
Jyothishmathi Institute of Technology \& Science, Karimnagar, during academic year
\textbf{2025--2026}, is a bonafide work carried out by them under my guidance and
supervision. The results presented in this Project Work have been verified and are
found to be satisfactory. The results embodied in this Project Work have not been
submitted to any other University for the award of any other degree or diploma.
}
\end{doublespace}

\vspace{3cm}

\noindent
\begin{tabular}{p{0.45\textwidth} p{0.1\textwidth} p{0.45\textwidth}}
  \textbf{Project Guide} & & \textbf{Head of the Department} \\[1.5cm]
  \textbf{K. Sai Veena} & & \textbf{Dr. M. Ravindar} \\
  Associate Professor   & & Associate Professor \\
  Dept. of CSE, JITS    & & Dept. of CSE, JITS \\
\end{tabular}

\vspace{2cm}
\begin{center}
  \textbf{EXTERNAL EXAMINER}
\end{center}

%============================================================
%  DECLARATION
%============================================================
\thispagestyle{empty}
\newpage
\begin{center}
  {\fontsize{16}{20}\selectfont\bfseries DECLARATION}
\end{center}
\vspace{2\baselineskip}

{\fontsize{12}{16}\selectfont
\onehalfspacing
We hereby declare that the work which is being presented in this dissertation entitled,
\textbf{``TALENT LINK: A PROFESSIONAL FREELANCE PLATFORM''}, submitted towards the
partial fulfilment of the requirements for the award of the degree of Bachelor of
Technology in Computer Science \& Engineering, Jyothishmathi Institute of Technology
\& Science, Karimnagar is an authentic record of our own work carried out under the
supervision of \textbf{K. Sai Veena}, Associate Professor, Department of CSE,
Jyothishmathi Institute of Technology and Science, Karimnagar.

To the best of our knowledge and belief, this project bears no resemblance with any
report submitted to JNTUH or any other University for the award of any degree or
diploma.
}

\vspace{3cm}
\noindent
\begin{tabular}{p{0.45\textwidth} p{0.1\textwidth} p{0.45\textwidth}}
  ADI ROSHINI (23271A05D9) & & \\[0.8cm]
  VATARIKARI MANISHA (23271A05C8) & & \\[0.8cm]
  ARMOOR AKHILA (23271A05B9) & & \\[0.8cm]
  ARANDKAR VAESHNAVI (23271A05G3) & & \\
\end{tabular}

\vspace{1.5cm}
\noindent Date: \hrulefill\\[0.5cm]
\noindent Place: Karimnagar

%============================================================
%  ACKNOWLEDGEMENT  (no page number)
%============================================================
\thispagestyle{empty}
\newpage
\begin{center}
  {\fontsize{16}{20}\selectfont\bfseries ACKNOWLEDGEMENT}
\end{center}
\vspace{2\baselineskip}
\onehalfspacing

{\fontsize{12}{16}\selectfont
\indent We would like to express our sincere gratitude to our advisor,
\textbf{Dr. M. RAVINDAR}, whose knowledge and guidance has motivated us to achieve
goals we never thought possible. The time we have spent working under his supervision
has truly been a pleasure.

\indent The experience from this kind of work is great and will be useful to us in the
future. We thank \textbf{Dr. M. RAVINDAR}, Associate Professor \& HOD, CSE Dept.,
for his effort, kind cooperation, guidance, and for encouraging us to do this work and
also for providing the facilities to carry out this work.

\indent It is a great pleasure to convey our thanks to \textbf{Dr. T. ANIL KUMAR},
Principal, Jyothishmathi Institute of Technology \& Science, and the College Management
for permitting us to undertake this project and providing excellent facilities to carry
out our project work.

\indent We thank all the Faculty members of the Department of Computer Science \&
Engineering for sharing their valuable knowledge with us. We extend our thanks to the
Technical Staff of the department for their valuable suggestions to technical problems.
Finally, special thanks to our parents for their support and encouragement throughout
our life and this course. Thanks to all our friends and well-wishers for their constant
support.
}

%============================================================
%  ABSTRACT  (page number starts here: i)
%============================================================
\newpage
\setcounter{page}{1}
\thispagestyle{plain}
\begin{center}
  {\fontsize{16}{20}\selectfont\bfseries ABSTRACT}
\end{center}
\vspace{2\baselineskip}
\addcontentsline{toc}{chapter}{Abstract}

{\fontsize{12}{16}\selectfont\onehalfspacing
\indent TalentLink connects freelancers and clients through a clear and organized system.
As remote work and the digital economy grow, there is a greater need for a reliable
platform. TalentLink is a freelance matchmaking platform that makes hiring and
collaboration easier. Current systems often struggle with disorganized workflows, gaps
in communication, and challenges in managing proposals and project progress. TalentLink
offers a dependable space where clients can post project needs, and freelancers can find
opportunities and submit proposals that match their skills.

\indent The platform promotes organized project management, clear communication, and
structured teamwork throughout the project. By streamlining workflows and improving user
communication, TalentLink makes managing freelance projects simpler. The main goal of the
system is to make the hiring process easier, improve transparency, and create a
professional environment for freelance collaboration. TalentLink seeks to boost
efficiency, support effective teamwork, and deliver a scalable solution.

\indent The rapid growth of digital technology and the rise of remote work have changed
how organizations and individuals work together. Freelancing platforms connect skilled
professionals with clients needing specific services for project-based work. However,
many existing freelance platforms struggle with issues like unstructured workflows, poor
project management, communication gaps, and challenges in finding reliable professionals.
These issues create difficulties for freelancers looking for real opportunities and for
clients seeking skilled individuals to complete projects effectively.

\indent TalentLink offers a structured workflow that supports proposal submission,
communication, project tracking, and collaboration throughout the project's lifecycle.
It helps freelancers highlight their skills and experience while enabling clients to
review proposals and pick the best candidates for their projects. By improving
communication and lowering workflow complexity, the system ensures smoother collaboration
and better project results.
}

%============================================================
%  TABLE OF CONTENTS
%============================================================
\newpage
\thispagestyle{plain}
{\setlength{\parskip}{0pt}
\tableofcontents}

%============================================================
%  LIST OF FIGURES
%============================================================
\newpage
\thispagestyle{plain}
{\setlength{\parskip}{0pt}
\listoffigures}
\addcontentsline{toc}{chapter}{List of Figures}

%============================================================
%  LIST OF TABLES
%============================================================
\newpage
\thispagestyle{plain}
{\setlength{\parskip}{0pt}
\listoftables}
\addcontentsline{toc}{chapter}{List of Tables}

%============================================================
%  LIST OF ABBREVIATIONS
%============================================================
\newpage
\thispagestyle{plain}
\begin{center}
  {\fontsize{16}{20}\selectfont\bfseries LIST OF ABBREVIATIONS}
\end{center}
\addcontentsline{toc}{chapter}{List of Abbreviations}
\vspace{2\baselineskip}

\begin{tabular}{p{0.2\textwidth} p{0.7\textwidth}}
  \textbf{Abbreviation} & \textbf{Full Form} \\[0.5cm]
  API      & Application Programming Interface \\
  AICTE    & All India Council for Technical Education \\
  CSE      & Computer Science and Engineering \\
  DFD      & Data Flow Diagram \\
  DRF      & Django REST Framework \\
  HOD      & Head of the Department \\
  HTML     & HyperText Markup Language \\
  HTTP     & HyperText Transfer Protocol \\
  IDE      & Integrated Development Environment \\
  JITS     & Jyothishmathi Institute of Technology and Science \\
  JNTUH    & Jawaharlal Nehru Technological University, Hyderabad \\
  MVP      & Minimum Viable Product \\
  NAAC     & National Assessment and Accreditation Council \\
  NBA      & National Board of Accreditation \\
  OTP      & One Time Password \\
  RAM      & Random Access Memory \\
  REST     & Representational State Transfer \\
  UI       & User Interface \\
  UML      & Unified Modelling Language \\
  URL      & Uniform Resource Locator \\
  VS Code  & Visual Studio Code \\
\end{tabular}

%============================================================
%  MAIN MATTER  (Arabic page numbering from 1)
%============================================================
\newpage
\pagenumbering{arabic}
\setcounter{page}{1}

%============================================================
%  CHAPTER 1 – INTRODUCTION
%============================================================
\chapter{INTRODUCTION}

\section{PROJECT OVERVIEW}

Freelancing has become a key part of the modern digital economy. It allows individuals
to work independently and offer services to clients from various locations. Many
organizations and businesses prefer hiring freelancers for short-term projects because of
the flexibility and access to specific skills this approach offers. However, finding
trustworthy freelancers and managing freelance projects can be tough for both clients and
freelancers. TalentLink is a professional freelance matchmaking platform designed to
connect freelancers and clients through a clear and organized system.

\indent The platform lets clients post project needs and allows freelancers to browse
available opportunities and submit proposals that match their skills. It also supports
project management, communication, and teamwork throughout the project. The main goal of
TalentLink is to make the freelance hiring process easier and improve collaboration
between freelancers and clients. By offering a centralized platform, the system ensures
transparency, an organized workflow, and effective management of freelance projects.

\indent To address the challenges faced in connecting talented individuals with suitable
opportunities, this project proposes TalentLink, an intelligent web-based platform
designed to bridge the gap between skilled individuals and organizations seeking talent.
The system aims to create a centralized and efficient environment where users can
showcase their abilities, portfolios, and achievements while enabling recruiters to
easily discover and connect with potential candidates. By integrating modern web
technologies and intelligent recommendation mechanisms, the platform ensures a smooth
and transparent talent discovery process.

\section{PROJECT PURPOSE}

The purpose of the TalentLink project is to create a professional platform that connects
freelancers and clients in a clear and effective environment. With the rise of the digital
economy and the remote work culture, freelancing has become a popular way for
organizations to complete projects using skilled professionals from various locations.
However, many freelancers struggle to find reliable project opportunities, while clients
often find it hard to identify the right talent for their needs. TalentLink seeks to solve
these problems by offering a centralized system that streamlines freelance hiring and
project management. The platform allows for smooth interaction between freelancers and
clients while promoting transparency and an organized workflow.

\indent TalentLink is created to simplify the entire freelance hiring lifecycle---from
project posting to proposal submission, selection, execution, and completion. The system
provides a streamlined approach where clients can clearly define their project
requirements, budgets, and timelines, while freelancers can showcase their skills,
experience, and previous work. By implementing a well-organized proposal management
system, the platform ensures that clients can easily compare submissions and select the
most suitable candidate based on expertise and credibility.

\indent TalentLink focuses on building a trustworthy and transparent freelance
environment. By maintaining professional profiles, organized records of transactions, and
secure interactions, the platform encourages accountability and long-term collaboration.
The system is designed not only to connect people but also to create a sustainable
freelance ecosystem where quality work, timely delivery, and client satisfaction are
prioritized. The platform aims to promote digital employment, support remote work
culture, and contribute to the growth of the freelance economy by offering a reliable,
scalable, and user-friendly solution.

\section{PROJECT SCOPE}

The TalentLink project aims to create a professional freelance matchmaking platform that
connects freelancers and clients in an efficient way. The platform simplifies finding
freelance opportunities and managing project collaborations. It offers a central system
where clients can post projects, and freelancers can browse available work and submit
proposals. TalentLink seeks to improve the freelance workflow by providing clear
communication, project tracking, and collaboration tools.

\indent The scope of the TalentLink project focuses on developing a professional
freelance matchmaking platform that connects freelancers and clients in an organized and
efficient environment. The key scope areas include:

\begin{itemize}[leftmargin=2em]
  \item \textbf{Project Posting and Proposal Management:} Clients can post their project
        needs, and freelancers can review the details and submit proposals based on their
        skills.
  \item \textbf{Freelancer and Client Interaction:} The platform allows freelancers and
        clients to communicate, making it easier to discuss project requirements and work
        together.
  \item \textbf{Project Tracking and Management:} TalentLink organizes project workflows
        so users can monitor progress from the proposal stage to completion.
  \item \textbf{User Profile Management:} Freelancers can highlight their skills and
        experience in their profiles, helping clients find the right candidates.
  \item \textbf{Scalability and Future Expansion:} The platform can grow to add features
        like smart job matching, better collaboration tools, and support for more users.
  \item \textbf{Secure User Authentication:} The system provides secure login and access
        control to ensure only authorized users access the platform.
  \item \textbf{Review and Feedback System:} Clients can provide ratings and feedback to
        freelancers after project completion, helping maintain trust and transparency.
  \item \textbf{Efficient Collaboration Environment:} The platform provides an organized
        environment that supports smooth communication and collaboration.
\end{itemize}

\section{PROJECT FEATURES}

\begin{itemize}[leftmargin=2em]
  \item \textbf{User Authentication \& Security:} TalentLink allows freelancers and
        clients to sign up for an account; once registered, users can log in with secure
        credentials. Only registered users can access the system.
  \item \textbf{Project Posting:} Clients may post project details (length of project,
        skill set, etc.) so freelancers can review and identify potential projects.
  \item \textbf{Proposal Submission:} Freelancers can review posted projects and submit
        proposals according to their skills, experience, and the type of work needed.
  \item \textbf{Communication:} The platform provides a means for freelancers and clients
        to communicate regarding project collaboration.
  \item \textbf{Project Management:} The platform provides users with a streamlined method
        for managing a project by tracking from proposal submission to completion.
  \item \textbf{Review and Feedback:} After a project is completed, clients can provide
        feedback and rate the freelancer's work, creating an environment of openness and
        trust.
  \item \textbf{User Profiles:} Freelancers can create a profile that lists their
        experience, skill set, and previous projects, helping clients find the right
        freelancer for the job.
\end{itemize}

%============================================================
%  CHAPTER 2 – LITERATURE REVIEW
%============================================================
\chapter{LITERATURE REVIEW}

The rapid rise in digital technology and remote work has led to an increase in demand for
online freelancing platforms. These platforms link clients that have a requirement for a
precise task or service with independent contractors who possess the skills needed to meet
those requirements. The way that freelancing systems work gives independent contractors
the ability to work as independent professionals. In addition, organizations benefit from
being able to access all of the global talent available through freelance systems.
However, having effective management of freelance workers and successful collaborations
between organizations and independent contractors requires an effective structure for how
projects are established, a way to have transparent communication among all parties, and a
reliable project management process.

\indent Numerous research documents exist regarding the development of online marketplace
platforms for service-based services between independent contractors and their clients.
These platforms are expected to provide users with features such as registration, project
posting, submitting proposals, and communication during execution. Developing and managing
the overall structure of the workflow in conjunction with secured communication between
independent contractors and their clients is critical to improving collaboration levels
and successfully completing the project.

\section{RELATED WORKS}

\subsection{Title 1: The Freelancer Application: A Gateway to Job-Seeking Opportunities for Students}

\noindent\textbf{Author:} Lily Puspa Dewi, Michael Wong, and Henry Novianus Palit\\
\textbf{Year:} 2022

The purpose of the study is to explore an E-commerce application designed to create a
marketplace for freelance jobs for college students. This platform serves as a place where
users are able to create an account and search for and apply to freelance positions they
believe their skills will be able to fulfil. Key components of this research include:
developing an application that allows users to use secure registration methods;
opportunities to develop a user workflow for finding quality applicants efficiently; and
creating an efficient method for communicating and completing transactions in a timely
manner, while improving client-to-freelancer collaboration and making the job search
experience easier.

\subsection{Title 2: Design of Online Service Marketplace Platforms}

\noindent\textbf{Author:} Geoffrey G. Parker, Marshall Van Alstyne, and Vili Lehdonvirta\\
\textbf{Year:} 2021

Research provides a framework for developing an online marketplace for services through
which businesses and customers can connect. Emphasising the need to establish structured
project management methodologies and be able to effectively evaluate proposals submitted
on the marketplace, allows users to identify the best candidates based on their areas of
expertise. Through this research project, centralised platforms improve the communication
and project management on the client side by improving the efficiency of finding an
approved method for completing the desired task.

\subsection{Title 3: Digital Freelance Marketplaces and Remote Work Platforms}

\noindent\textbf{Author:} Vili Lehdonvirta, John Horton, and Arun Sundararajan\\
\textbf{Year:} 2020

The research looks into how digital freelancing marketplaces are growing and how they are
affecting global labour markets. The research explains that businesses can get access to
specialised skills via freelancing platforms and people gain the benefits of flexible
work. The research also highlights the need for transparency, trust, and structured
collaboration tools for successful freelance project management. From the analysis of
these research studies, TalentLink has been built as a professional matchmaking platform
for freelancers, providing a structured workflow, organised communication, and efficient
project management.

%============================================================
%  CHAPTER 3 – EXISTING AND PROPOSED SYSTEM
%============================================================
\chapter{EXISTING AND PROPOSED SYSTEM}

\section{EXISTING SYSTEM}

The existing freelancing systems provide platforms where clients can post projects, and
freelancers can search for job opportunities. These platforms allow freelancers to apply
for projects based on their skills and experience. However, many existing systems still
face challenges in managing freelance collaborations efficiently. In several cases, the
workflow between freelancers and clients is not well organized, which may lead to
confusion and delays in project completion.

\subsection{Existing System Disadvantages}

\begin{itemize}[leftmargin=2em]
  \item Lack of structured workflow for managing freelance projects.
  \item Difficulty in identifying reliable freelancers for specific project requirements.
  \item Communication gaps between freelancers and clients.
  \item Inefficient project tracking and management.
  \item Time-consuming hiring process.
\end{itemize}

\section{PROBLEM STATEMENT}

In the current freelance ecosystem, clients often face difficulties in finding skilled and
reliable freelancers who can meet their project requirements within the specified time.
Similarly, freelancers may struggle to discover genuine and suitable project opportunities
that match their skills and expertise. Many existing freelance systems do not provide a
structured workflow for managing freelance collaborations, which may lead to communication
gaps, delays in project completion, and inefficiencies in project management.

\indent Additionally, the lack of a centralized platform for handling project proposals,
monitoring project progress, and maintaining clear communication between users makes the
freelance hiring process more complicated and time-consuming.

\indent Therefore, there is a need to develop a professional freelance matchmaking
platform that simplifies the hiring process and improves collaboration between freelancers
and clients. The system should:

\begin{itemize}[leftmargin=2em]
  \item Provide a centralized platform where clients can post project requirements.
  \item Allow freelancers to browse available projects and submit proposals based on their
        skills.
  \item Enable effective communication between freelancers and clients.
  \item Support organized project tracking and management.
  \item Improve transparency and efficiency in freelance collaboration.
\end{itemize}

\section{PROPOSED SYSTEM}

The limitations of existing freelance systems highlight the need for a more efficient
platform for managing freelance collaborations. TalentLink is proposed as a professional
freelance matchmaking platform that provides a structured environment for connecting
freelancers and clients. The system allows clients to post project requirements and
enables freelancers to browse available opportunities and submit proposals based on their
skills and experience.

\indent The platform ensures organized communication and collaboration between freelancers
and clients throughout the project lifecycle. By providing a centralized system for
project posting, proposal submission, and project tracking, TalentLink simplifies the
freelance hiring process and improves transparency. The system helps clients identify
suitable freelancers efficiently while enabling freelancers to access genuine project
opportunities.

\indent TalentLink also supports structured project management, allowing users to monitor
project progress and maintain clear communication. By improving workflow organization and
collaboration, the platform aims to create a reliable and professional environment for
freelance project management.

\subsection{Proposed System Advantages}

\begin{itemize}[leftmargin=2em]
  \item \textbf{Efficient Hiring Process:} The platform simplifies the process of finding
        and hiring freelancers by providing a centralized system for project posting and
        proposal submission.
  \item \textbf{Improved Communication:} TalentLink enables clear and organized
        communication between freelancers and clients, reducing misunderstandings during
        project collaboration.
  \item \textbf{Better Project Management:} The system provides structured workflows for
        managing projects, making it easier to track progress and monitor activities.
  \item \textbf{Transparency and Trust:} The platform allows clients to review freelancer
        profiles and feedback, helping build trust and transparency in freelance
        collaborations.
  \item \textbf{Access to Genuine Opportunities:} Freelancers can easily discover relevant
        project opportunities that match their skills and expertise.
  \item \textbf{Centralized Platform:} All freelance activities such as project posting,
        proposal management, and communication are handled within a single platform,
        improving efficiency.
\end{itemize}

%============================================================
%  CHAPTER 4 – SYSTEM REQUIREMENTS
%============================================================
\chapter{SYSTEM REQUIREMENTS}

\section{HARDWARE REQUIREMENTS}

Hardware requirements specify the basic physical components required to develop and run
the TalentLink platform. Since the project is a web-based application, it does not require
any specialized hardware components. The following hardware components are sufficient for
the development and execution of the system:

\begin{table}[H]
  \centering
  \caption{Hardware Requirements}
  \begin{tabularx}{0.85\textwidth}{|>{\bfseries}l|X|}
    \hline
    \textbf{Component}       & \textbf{Specification} \\
    \hline
    Processor                & Intel Core i3 or above \\
    \hline
    RAM                      & Minimum 4 GB (8 GB recommended) \\
    \hline
    Storage                  & Minimum 10 GB free disk space \\
    \hline
    System Type              & 64-bit Computer \\
    \hline
    Internet Connection      & Required for development, testing, and deployment \\
    \hline
  \end{tabularx}
  \label{tab:hardware}
\end{table}

\section{SOFTWARE REQUIREMENTS}

Software requirements define the tools and technologies used to develop the TalentLink
platform.

\begin{table}[H]
  \centering
  \caption{Software Requirements}
  \begin{tabularx}{0.85\textwidth}{|>{\bfseries}l|X|}
    \hline
    \textbf{Category}        & \textbf{Tool / Technology} \\
    \hline
    Operating System         & Windows / Linux / macOS \\
    \hline
    Programming Languages    & Python, JavaScript \\
    \hline
    Frontend Framework       & React \\
    \hline
    Backend Framework        & Django, Django REST Framework \\
    \hline
    Database                 & SQLite (Development), PostgreSQL (Production) \\
    \hline
    API Tools                & Postman, Swagger \\
    \hline
    IDE                      & Visual Studio Code \\
    \hline
    Version Control          & Git, GitHub \\
    \hline
    Web Browser              & Google Chrome \\
    \hline
  \end{tabularx}
  \label{tab:software}
\end{table}

%============================================================
%  CHAPTER 5 – PROJECT DESCRIPTION
%============================================================
\chapter{PROJECT DESCRIPTION}

The TalentLink platform is designed to provide a professional environment where
freelancers and clients can collaborate efficiently for project-based work. The system
focuses on simplifying the process of discovering freelance opportunities, submitting
proposals, and managing projects through a structured workflow.

\indent The platform consists of two primary user roles: Clients and Freelancers. Clients
are users who post project requirements, while freelancers are professionals who search
for projects and submit proposals based on their skills and experience. The system enables
smooth interaction between these two roles through organized communication and project
management features.

\indent When a client registers on the platform, they can create and publish project
listings that include details such as project description, required skills, deadlines, and
budget. Freelancers can browse available projects and submit proposals that describe how
they can complete the project. Clients can review these proposals and select the most
suitable freelancer for their project.

\indent Once a freelancer is selected, the platform allows both users to communicate and
collaborate throughout the project lifecycle. The system supports project tracking,
enabling both freelancers and clients to monitor the progress of the project. This
structured workflow helps ensure transparency and reduces misunderstandings between users.

\indent TalentLink also provides profile management features that allow freelancers to
showcase their skills, experience, and previous work. This helps clients evaluate
freelancers effectively before assigning projects. Similarly, clients can manage their
projects, review proposals, and track project activities through the platform.

\section{SETTING UP THE TALENTLINK APPLICATION}

\subsection{Step 1: User Registration}

Users first register on the TalentLink platform by providing basic details such as name,
email, and password. The system verifies the user using an OTP-based authentication
process to ensure secure access.

\begin{figure}[H]
  \centering
  % \includegraphics[width=0.75\textwidth]{signup.png}
  \fbox{\parbox{0.75\textwidth}{\centering\vspace{2cm}[Fig 5.1: TalentLink Sign Up Screenshot]\vspace{2cm}}}
  \caption{TalentLink Sign Up}
  \label{fig:signup}
\end{figure}

\subsection{Step 2: User Login}

After successful registration, users can log in to the system using their credentials.
Based on their role, they are directed to either the freelancer dashboard or the client
dashboard.

\begin{figure}[H]
  \centering
  \fbox{\parbox{0.75\textwidth}{\centering\vspace{2cm}[Fig 5.2: Client Dashboard Screenshot]\vspace{2cm}}}
  \caption{TalentLink Client Dashboard}
  \label{fig:client_dash}
\end{figure}

\begin{figure}[H]
  \centering
  \fbox{\parbox{0.75\textwidth}{\centering\vspace{2cm}[Fig 5.3: Freelancer Dashboard Screenshot]\vspace{2cm}}}
  \caption{TalentLink Freelancer Dashboard}
  \label{fig:freelancer_dash}
\end{figure}

\begin{figure}[H]
  \centering
  \fbox{\parbox{0.75\textwidth}{\centering\vspace{2cm}[Fig 5.4: Login Screen Screenshot]\vspace{2cm}}}
  \caption{TalentLink Login}
  \label{fig:login}
\end{figure}

\subsection{Step 3: Profile Creation}

Freelancers can create a detailed profile that lists their skills, experience, and
previous work so clients can evaluate them effectively.

\begin{figure}[H]
  \centering
  \fbox{\parbox{0.75\textwidth}{\centering\vspace{2cm}[Fig 5.5: Profile Creation Screenshot]\vspace{2cm}}}
  \caption{Profile Creation}
  \label{fig:profile}
\end{figure}

\subsection{Step 4: Project Posting}

Clients can create and publish project listings with details such as project description,
required skills, and deadlines.

\begin{figure}[H]
  \centering
  \fbox{\parbox{0.75\textwidth}{\centering\vspace{2cm}[Fig 5.6: Post Project Screenshot]\vspace{2cm}}}
  \caption{Post Project -- Client View}
  \label{fig:post_project1}
\end{figure}

\begin{figure}[H]
  \centering
  \fbox{\parbox{0.75\textwidth}{\centering\vspace{2cm}[Fig 5.7: Post Project Form Screenshot]\vspace{2cm}}}
  \caption{Post Project -- Project Form}
  \label{fig:post_project2}
\end{figure}

\subsection{Step 5: Proposal Submission}

Freelancers can browse available projects and submit proposals explaining how they plan
to complete the project.

\begin{figure}[H]
  \centering
  \fbox{\parbox{0.75\textwidth}{\centering\vspace{2cm}[Fig 5.8: Proposal Submission Screenshot]\vspace{2cm}}}
  \caption{Proposal Submission}
  \label{fig:proposal}
\end{figure}

\subsection{Step 6: Project Selection}

Clients review proposals submitted by freelancers and select the most suitable candidate
for their project.

\begin{figure}[H]
  \centering
  \fbox{\parbox{0.75\textwidth}{\centering\vspace{2cm}[Fig 5.9: Project Selection Screenshot]\vspace{2cm}}}
  \caption{Project Selection}
  \label{fig:selection}
\end{figure}

\subsection{Step 7: Communication and Collaboration}

Once a freelancer is selected, both parties can communicate through the platform and work
on the project collaboratively.

\begin{figure}[H]
  \centering
  \fbox{\parbox{0.75\textwidth}{\centering\vspace{2cm}[Fig 5.10: Messaging Interface Screenshot]\vspace{2cm}}}
  \caption{Messaging and Collaboration}
  \label{fig:messaging}
\end{figure}

\subsection{Step 8: Project Completion and Review}

After completing the project, clients can provide feedback and ratings to freelancers,
helping maintain transparency and trust in the platform.

%============================================================
%  CHAPTER 6 – SYSTEM DESIGN
%============================================================
\chapter{SYSTEM DESIGN}

\section{SYSTEM ARCHITECTURE}

The system architecture of TalentLink follows a structured web application architecture
that enables communication between users, the application interface, and the database.
The platform is designed to provide a reliable environment where freelancers and clients
can interact, manage projects, and collaborate efficiently.

\indent The architecture consists of three main components: the frontend layer, the
backend layer, and the database layer. The frontend layer provides the user interface
through which clients and freelancers interact with the system. It allows users to
register, log in, browse projects, submit proposals, and manage project activities.

\indent The database layer is responsible for storing and managing all system data,
including user information, project details, proposals, contracts, and reviews. The
database ensures that information is stored securely and can be retrieved whenever
required.

\indent When a user performs an action on the platform, the request is sent from the
frontend to the backend through application interfaces. The backend processes the
request, interacts with the database to store or retrieve data, and then sends the
response back to the frontend. This structured architecture helps maintain system
efficiency, scalability, and organized management of freelance project activities.

\section{USER FLOW DESIGN}

A Data Flow Diagram (DFD) is a graphical representation of the flow of data within a
system or process. It is a modelling technique that shows how data moves through
processes and how it is stored, transformed, or exchanged within a system. DFDs are often
used in system analysis and design to visualize and understand the data processing and
flow of information in a structured manner.

\begin{figure}[H]
  \centering
  \fbox{\parbox{0.85\textwidth}{\centering\vspace{3cm}[Fig 6.1: User Flow Diagram]\vspace{3cm}}}
  \caption{User Flow Diagram for TalentLink}
  \label{fig:userflow}
\end{figure}

\section{UML DIAGRAMS}

Unified Modelling Language (UML) is a standardized modelling language in the field of
software engineering. It provides a way to visualize a system's design through a set of
diagrams. UML diagrams help software developers and system architects communicate and
understand the structure and behaviour of a system. In UML, the diagrams can be broadly
categorized into two main types: structural diagrams and behavioural diagrams.

\subsection{Structural Diagrams}

Structural diagrams in UML focus on representing the static structure of a system. They
showcase the components that make up the system and their relationships. The following
are different structural diagrams:

\begin{itemize}[leftmargin=2em]
  \item Class Diagram
  \item Object Diagram
  \item Component Diagram
  \item Deployment Diagram
\end{itemize}

\subsection{Behavioural Diagrams}

Behavioural diagrams primarily capture the dynamic facet of the system. UML has the
following 5 varieties of behavioural diagrams:

\begin{itemize}[leftmargin=2em]
  \item Use Case Diagram
  \item Sequence Diagram
  \item Collaboration Diagram
  \item State Chart Diagram
  \item Activity Diagram
\end{itemize}

\subsubsection{6.3.1 Use Case Diagram}

Use Case Diagrams in UML describe interactions between a system and external entities
known as actors. Use cases represent specific functionalities or scenarios that the
system provides to its users.

\begin{figure}[H]
  \centering
  \fbox{\parbox{0.85\textwidth}{\centering\vspace{3cm}[Fig 6.3.1: Use Case Diagram]\vspace{3cm}}}
  \caption{Use Case Diagram}
  \label{fig:usecase}
\end{figure}

\subsubsection{6.3.2 Class Diagram}

The Class Diagram in UML illustrates the static structure of a system, detailing classes,
attributes, methods, and their relationships. Rectangles represent classes, and lines
indicate associations, dependencies, and inheritances.

\begin{figure}[H]
  \centering
  \fbox{\parbox{0.85\textwidth}{\centering\vspace{3cm}[Fig 6.3.2: Class Diagram]\vspace{3cm}}}
  \caption{Class Diagram}
  \label{fig:class}
\end{figure}

\subsubsection{6.3.3 Activity Diagram}

Activity Diagrams represent the flow of activities within a process or workflow. They
focus on actions, decisions, and control flows, providing a high-level view of the
dynamic aspects of a system.

\begin{figure}[H]
  \centering
  \fbox{\parbox{0.85\textwidth}{\centering\vspace{3cm}[Fig 6.3.3: Activity Diagram]\vspace{3cm}}}
  \caption{Activity Diagram}
  \label{fig:activity}
\end{figure}

\subsubsection{6.3.4 Sequence Diagram}

A sequence diagram depicts interaction between objects in a serial order, i.e., the
order during which these interactions occur. Sequence diagrams are widely employed by
business analysts and software developers to document and understand requirements for new
and existing systems.

\begin{figure}[H]
  \centering
  \fbox{\parbox{0.85\textwidth}{\centering\vspace{3cm}[Fig 6.3.4: Sequence Diagram]\vspace{3cm}}}
  \caption{Sequence Diagram}
  \label{fig:sequence}
\end{figure}

%============================================================
%  CHAPTER 7 – IMPLEMENTATION
%============================================================
\chapter{IMPLEMENTATION}

The TalentLink platform was implemented as a full-stack web application designed to
connect freelancers and clients through a structured and user-friendly interface. The
system allows users to register and authenticate securely, after which they are assigned
roles as either freelancers or clients. Clients can post projects by providing details
such as title, description, budget, and required skills. Freelancers can browse available
projects and submit proposals based on their expertise. The platform also enables clients
to review proposals and select suitable freelancers for their projects, ensuring an
organized hiring process.

\indent The backend of the system manages user authentication, project management,
proposal handling, messaging, and contract tracking. All requests from the frontend are
processed through secure APIs and stored in the database for efficient retrieval and
management. The system also supports communication between freelancers and clients and
allows users to track project progress from proposal submission to completion. This
structured implementation ensures transparency, efficient collaboration, and a smooth
workflow for freelance project management.

\section{INSTALLATIONS AND LIBRARIES USED}

\subsection{Installation Steps}

\begin{enumerate}[leftmargin=2em]
  \item Install Python and create a virtual environment.
  \item Install backend dependencies using \texttt{pip} from the requirements file.
  \item Install Node.js and npm for the frontend development environment.
  \item Install frontend dependencies using \texttt{npm install}.
  \item Configure the database connection for SQLite (development) or PostgreSQL
        (production).
  \item Run database migrations and start the backend server.
  \item Start the frontend development server to launch the application.
\end{enumerate}

\section{CODE}

The following is the main routing file (\texttt{App.js}) of the React frontend, which
defines all application routes for both the Client and Freelancer portals.

\begin{verbatim}
import React from "react";
import "./App.css";
import "./assets/global.css";
import { BrowserRouter as Router, Routes, Route, Navigate }
  from "react-router-dom";

/* ===== Context Providers ===== */
import { UserProvider } from "./context/UserContext";
import { ProjectProvider } from "./context/ProjectContext";
import { ThemeProvider } from "./context/ThemeContext";
import { SearchProvider } from "./context/SearchContext";

/* ===== AuthFlow UI Pages ===== */
import Home from "./pages/Home";
import Login from "./pages/Login";
import Signup from "./pages/Signup";
import VerifyOtp from "./pages/VerifyOtp";
import ForgotPassword from "./pages/ForgotPassword";
import ResetPassword from "./pages/ResetPassword";

/* ===== Client Dashboard Pages & Layout ===== */
import ClientLayout from "./freelancer_layouts/ClientLayout";
import ClientDashboard from "./freelancer_pages/client/ClientDashboard";
import ClientProjects from "./freelancer_pages/client/ClientProjects";

/* ===== Freelancer Dashboard Pages ===== */
import FreelancerDashboard
  from "./freelancer_pages/freelancer/FreelancerDashboard";
import Projects from "./freelancer_pages/freelancer/Projects";
import FreelancerLayout
  from "./freelancer_layouts/FreelancerLayout";

function App() {
  return (
    <ThemeProvider>
      <UserProvider>
        <ProjectProvider>
          <SearchProvider>
            <Router>
              <Routes>
                {/* PUBLIC ROUTES */}
                <Route path="/" element={<Home />} />
                <Route path="/login" element={<Login />} />
                <Route path="/signup" element={<Signup />} />
                <Route path="/verify-otp" element={<VerifyOtp />} />

                {/* CLIENT ROUTES */}
                <Route path="/client" element={<ClientLayout />}>
                  <Route index element={<ClientDashboard />} />
                  <Route path="projects" element={<ClientProjects />} />
                </Route>

                {/* FREELANCER ROUTES */}
                <Route path="/freelancer" element={<FreelancerLayout />}>
                  <Route index element={<FreelancerDashboard />} />
                  <Route path="projects" element={<Projects />} />
                </Route>
              </Routes>
            </Router>
          </SearchProvider>
        </ProjectProvider>
      </UserProvider>
    </ThemeProvider>
  );
}

export default App;
\end{verbatim}

%============================================================
%  CHAPTER 8 – RESULTS
%============================================================
\chapter{RESULTS}

This section describes step-by-step the system analysis and the results obtained. The
TalentLink platform was tested to verify the correct functioning of each module including
user authentication, project posting, proposal submission, messaging, and project
tracking. All functional modules were verified and found to be working as expected.

\indent The system provides a seamless user experience across both the Client and
Freelancer dashboards. Users can register, log in securely with OTP verification, post
projects, submit and review proposals, communicate through messaging, and track project
progress. The interface is responsive and intuitive, supporting smooth navigation across
all platform features.

\indent The backend APIs were tested using Postman, and all endpoints responded correctly
with appropriate status codes and data formats. The database correctly stored and
retrieved all information including user data, project details, proposals, and review
records.

\begin{figure}[H]
  \centering
  \fbox{\parbox{0.85\textwidth}{\centering\vspace{3cm}[Fig 8.1: System Output Screenshot]\vspace{3cm}}}
  \caption{System Output -- TalentLink Platform}
  \label{fig:output}
\end{figure}

%============================================================
%  CHAPTER 9 – FUTURE SCOPE
%============================================================
\chapter{FUTURE SCOPE}

While the TalentLink platform successfully provides a structured environment for
connecting freelancers and clients, there are several opportunities for further
improvements and enhancements. As the demand for digital freelance platforms continues to
grow, the system can be expanded with additional features to improve efficiency, user
experience, and scalability.

\indent The future scope of the TalentLink platform includes the following enhancements:

\begin{itemize}[leftmargin=2em]
  \item \textbf{AI-Based Project Matching:} Implement intelligent algorithms that
        automatically match freelancers with suitable projects based on skills,
        experience, and project requirements.
  \item \textbf{Advanced Communication Tools:} Integrate real-time chat, video
        communication, and file sharing features to improve collaboration between
        freelancers and clients.
  \item \textbf{Secure Payment Integration:} Add secure online payment systems to manage
        project payments, milestones, and financial transactions within the platform.
  \item \textbf{Mobile Application Development:} Develop a mobile version of the
        TalentLink platform to allow users to access freelance opportunities and manage
        projects from anywhere.
  \item \textbf{Enhanced Rating and Review System:} Improve the feedback system to help
        clients evaluate freelancers more effectively and maintain trust within the
        platform.
  \item \textbf{Blockchain-Based Contracts:} Introduce smart contracts using blockchain
        technology to ensure secure, tamper-proof agreements between freelancers and
        clients.
  \item \textbf{Multi-language Support:} Add support for multiple languages to make the
        platform accessible to a wider global audience.
\end{itemize}

%============================================================
%  CHAPTER 10 – CONCLUSION
%============================================================
\chapter{CONCLUSION}

The TalentLink platform provides a structured and efficient solution for connecting
freelancers and clients in a professional environment. With the increasing demand for
remote work and freelance services, there is a need for reliable platforms that simplify
the process of hiring skilled professionals and managing project collaborations.
TalentLink addresses these needs by providing a centralized system where clients can post
project requirements and freelancers can discover and apply for suitable opportunities.

\indent The system improves collaboration by enabling organized communication, proposal
submission, and project tracking throughout the project lifecycle. By offering a
structured workflow, TalentLink helps reduce communication gaps and improves transparency
between freelancers and clients. The platform also allows freelancers to showcase their
skills and experience, helping clients identify suitable professionals for their projects
more efficiently.

\indent Overall, TalentLink contributes to improving the freelance hiring process by
providing a reliable and organized platform for project collaboration. The system enhances
efficiency, transparency, and communication in freelance project management. In the
future, the platform can be further enhanced with advanced features such as intelligent
project matching, improved collaboration tools, and scalability to support a larger number
of users in the growing digital economy.

%============================================================
%  REFERENCES / BIBLIOGRAPHY
%============================================================
\newpage
\begin{center}
  {\fontsize{16}{20}\selectfont\bfseries REFERENCES}
\end{center}
\addcontentsline{toc}{chapter}{References}
\vspace{2\baselineskip}

\begin{enumerate}[leftmargin=2em, label={[\arabic*]}]

  \item Lily Puspa Dewi, Michael Wong, and Henry Novianus Palit,
        \textit{``The Freelancer Application: A Gateway to Job-Seeking Opportunities for
        Students,''} International Journal of Advanced Computer Science and Applications,
        2022.

  \item Geoffrey G. Parker, Marshall Van Alstyne, and Vili Lehdonvirta,
        \textit{``Design of Online Service Marketplace Platforms,''} Journal of
        Management Information Systems, 2021.

  \item Vili Lehdonvirta, John Horton, and Arun Sundararajan,
        \textit{``Digital Freelance Marketplaces and Remote Work Platforms,''} Oxford
        Internet Institute, University of Oxford, 2020.

  \item Django Software Foundation, \textit{Django Documentation},
        \url{https://docs.djangoproject.com}, 2024.

  \item React Documentation, \textit{React -- A JavaScript Library for Building User
        Interfaces}, \url{https://reactjs.org/docs/getting-started.html}, 2024.

  \item PostgreSQL Global Development Group, \textit{PostgreSQL Documentation},
        \url{https://www.postgresql.org/docs/}, 2024.

  \item GitHub Documentation, \textit{Git and GitHub Version Control},
        \url{https://docs.github.com}, 2024.

  \item Postman Inc., \textit{Postman API Testing Documentation},
        \url{https://learning.postman.com/docs/}, 2024.

\end{enumerate}

%============================================================
%  APPENDIX
%============================================================
\appendix
\chapter{APPENDIX}
\addcontentsline{toc}{chapter}{Appendix}

\section{TECHNOLOGY STACK SUMMARY}

\begin{table}[H]
  \centering
  \caption{Complete Technology Stack}
  \begin{tabularx}{0.85\textwidth}{|l|X|}
    \hline
    \textbf{Layer}       & \textbf{Technology Used}     \\
    \hline
    Frontend             & React.js, HTML5, CSS3        \\
    \hline
    Backend              & Django, Django REST Framework \\
    \hline
    Database             & SQLite / PostgreSQL           \\
    \hline
    Authentication       & OTP-based, JWT Tokens         \\
    \hline
    API Testing          & Postman, Swagger              \\
    \hline
    Version Control      & Git, GitHub                  \\
    \hline
    Deployment           & Local / Cloud                \\
    \hline
  \end{tabularx}
  \label{tab:techstack}
\end{table}

\section{PROJECT TEAM}

\begin{table}[H]
  \centering
  \caption{Project Team Details}
  \begin{tabularx}{0.85\textwidth}{|l|l|X|}
    \hline
    \textbf{S.No} & \textbf{Roll No.} & \textbf{Name} \\
    \hline
    1 & 23271A05D9 & ADI ROSHINI \\
    \hline
    2 & 23271A05C8 & VATARIKARI MANISHA \\
    \hline
    3 & 23271A05B9 & ARMOOR AKHILA \\
    \hline
    4 & 23271A05G3 & ARANDKAR VAESHNAVI \\
    \hline
  \end{tabularx}
  \label{tab:team}
\end{table}

%============================================================
\end{document}
%============================================================
